% datastruct/datastruct.tex
% mainfile: ../perfbook.tex
% SPDX-License-Identifier: CC-BY-SA-3.0

\QuickQuizChapter{chp:Data Structures}{数据结构}{qqzdatastruct}
%
\Epigraph{差劲的程序员操心代码。
	  优秀的程序员操心数据结构及其相互关系。}
	 {Linus Torvalds}

对算法的严肃讨论都会涉及其数据结构的时间复杂度~\cite{ThomasHCorman2001Algorithms}。
然而，对于并行程序，时间复杂度还包括并发效应，因为这些效应可能非常巨大，如
\cref{chp:Hardware and its Habits}所示。
换句话说，优秀程序员关注的数据结构关系中也包括与并发相关的方面。

本章将揭示若干复杂情况：

\begin{enumerate}
\item	完全按照\cref{chp:Partitioning and Synchronization Design}中的良好建议
	设计的数据结构，在某些类型的系统上仍然可能完全无法扩展。
\label{sec:datastruct:NUMA-induced failure to scale}
\item	完全按照\cref{chp:Partitioning and Synchronization Design}和
	\cref{chp:Deferred Processing}两章中的良好建议设计的数据结构，
	在某些类型的系统上\emph{仍然}可能完全无法扩展。
\label{sec:datastruct:Cache-size-induced failure to scale}
\item	即使是\emph{只读且无需同步}的数据结构遍历，
	在某些类型的系统上也可能无法扩展。
\label{sec:datastruct:Cache-size-induced failure to scale redux}
\item	避免了上述复杂情况的数据结构遍历，仍然可能受到并发更新的阻碍。
\label{sec:datastruct:Update-activity failure to scale}
\end{enumerate}

本章将研究这些复杂情况，并展示解决它们的一些方法。

\Cref{sec:datastruct:Motivating Application}
介绍本章数据结构的应用背景。
\Cref{chp:Partitioning and Synchronization Design}展示了分区如何提高可扩展性，因此
\cref{sec:datastruct:Partitionable Data Structures}
讨论可分区的数据结构。
\Cref{chp:Deferred Processing}描述了延迟某些操作如何能大幅提高性能和可扩展性，
\cref{sec:datastruct:Read-Mostly Data Structures}将继续探讨这一主题。
\Cref{sec:datastruct:Non-Partitionable Data Structures}
研究一种不可分区的数据结构，将其拆分为以读为主的部分和可分区的部分，
从而提高性能和可扩展性。
由于本章无法深入讨论每一种并发数据结构，
\cref{sec:datastruct:Other Data Structures}
将概述几种重要的数据结构。
最后，\cref{sec:datastruct:Summary}
对本章进行总结。

\section{激励应用}
\label{sec:datastruct:Motivating Application}
%
\epigraph{数学的艺术在于找到那个包含一切普遍性萌芽的特殊情形。}
	 {David Hilbert}

我们将使用 Schr\"odinger 的动物园应用来评估
性能~\cite{McKenney:2013:SDS:2483852.2483867}。
Schr\"odinger 拥有一个包含大量动物的动物园，
他希望使用内存数据库来追踪这些动物，数据库中的每个数据项代表动物园中的一只动物。
每只动物都有一个唯一的名字作为 key，并为每只动物跟踪各种数据。

出生、捕获和购买导致插入操作，而死亡、释放和出售导致删除操作。
因为 Schr\"odinger 的动物园包含大量短命的动物（包括老鼠和昆虫），
数据库必须处理高更新速率。
对 Schr\"odinger 动物感兴趣的人可以查询它们，
而 Schr\"odinger 注意到对他的猫的查询率高得可疑，
以至于他怀疑是他的老鼠在监视它们的天敌。
无论查询来自何处，Schr\"odinger 的应用都必须能处理对单个数据元素的高查询率。

我们将看到，这个简单的应用可能对传统的并发数据结构构成挑战。

\section{可分区数据结构}
\label{sec:datastruct:Partitionable Data Structures}
%
\epigraph{在当今找到一种过简单生活的方式，是最复杂的任务。}
	 {Henry A. Courtney，改编}

当今使用的数据结构种类繁多，多到有多本教科书专门介绍它们。
本节聚焦于一种数据结构，即 hash table。
这种聚焦的方法使我们能够更深入地研究并发如何与数据结构交互，
同时也关注了一种在实践中被大量使用的数据结构。
\Cref{sec:datastruct:Hash-Table Design}
概述设计，
\cref{sec:datastruct:Hash-Table Implementation}
给出实现。
最后，
\cref{sec:datastruct:Hash-Table Performance}
讨论由此产生的性能和可扩展性。

\subsection{哈希表（hash table）设计}
\label{sec:datastruct:Hash-Table Design}

\Cref{chp:Partitioning and Synchronization Design}
强调了需要应用分区来获得可观的性能和可扩展性，因此可分区性
必须是选择数据结构时的首要标准。
这一标准被并行计算的主力——hash table——很好地满足了。
Hash table 在概念上很简单，由一个\emph{hash bucket}数组组成。
\emph{Hash 函数}将给定元素的\emph{key}映射到该元素将被存储的 hash bucket。
因此每个 hash bucket 都引导着一个元素的 linked list，
称为\emph{hash chain}。
配置得当时，这些 hash chain 会非常短，
使得 hash table 能够极其高效地访问其元素。

\QuickQuiz{
	但是链式 hash table 只是众多类型中的一种。
	为什么要聚焦于链式 hash table？
}\QuickQuizAnswer{
	链式 hash table 是完全可分区的，因此非常适合并发使用。
	还有其他完全可分区的 hash table，
	例如 split-ordered list~\cite{OriShalev2006SplitOrderListHash}，
	但它们要复杂得多。
	因此我们聚焦于链式 hash table。
}\QuickQuizEnd

此外，每个 bucket 都有自己的 lock，因此 hash table 中不同 bucket 的元素可以
完全独立地进行添加、删除和查找。
一个拥有大量 bucket（因此也有大量 lock）的大型 hash table，
每个 bucket 包含少量元素，因此应该能提供优秀的可扩展性。

这种方法解决了\cref{sec:intro:Existing Sequential Designs}中指出的元素计数问题，
方法是不让 \co{hashtab_add()} 或 \co{hashtab_del()} 返回 hash table 中的元素数量。
它还通过允许重复元素以及将指向待删除元素的指针传递给 \co{hashtab_del()} 来进一步
简化 API。
这两项更改消除了返回错误的可能性，
这意味着 \co{hashtab_add()} 和 \co{hashtab_del()} 函数的返回类型可以是 \co{void}。

\subsection{哈希表实现}
\label{sec:datastruct:Hash-Table Implementation}

\begin{fcvref}[ln:datastruct:hash_bkt:struct]
\Cref{lst:datastruct:Hash-Table Data Structures}
(\path{hash_bkt.c})
展示了一组用于简单固定大小 hash table 的数据结构，
该 hash table 使用链接法和每 hash bucket 加锁的方式，
\cref{fig:datastruct:Hash-Table Data-Structure Diagram}
用图示说明了它们如何组合在一起。
注意，\co{cds_} 函数和数据结构可以在用户空间 RCU
库中找到~\cite{MathieuDesnoyers2009URCU,PaulMcKenney2013LWNURCUqueuestack,PaulMcKenney2013LWNURCUqueuestackAPI,PaulMcKenney2013LWNURCUlist}。
\co{hashtab} 结构（\cref{lst:datastruct:Hash-Table Data Structures}中的
\clnrefrange{tab:b}{tab:e}）
包含四个 \co{ht_bucket} 结构
（\cref{lst:datastruct:Hash-Table Data Structures}中的
\clnrefrange{bucket:b}{bucket:e}），
其中 \co{->ht_nbuckets} 字段控制 bucket 数量，
\co{->ht_cmp} 字段保存指向 key 比较函数的指针。
每个这样的 bucket 包含一个链表头 \co{->htb_head} 和
一个 lock \co{->htb_lock}。
链表头通过 \co{->hte_next} 字段将 \co{ht_elem} 结构
（\cref{lst:datastruct:Hash-Table Data Structures}中的
\clnrefrange{elem:b}{elem:e}）
串联起来，每个 \co{ht_elem} 结构还在 \co{->hte_hash} 字段中
缓存了对应元素的 hash 值。
\co{ht_elem} 结构被包含在一个更大的结构中，
该结构可能包含复杂的 key。
\end{fcvref}
\Cref{fig:datastruct:Hash-Table Data-Structure Diagram}
显示 bucket~0 包含两个元素，bucket~2 包含一个元素。

\begin{listing}
\input{CodeSamples/datastruct/hash/hash_bkt@struct.fcv}
\caption{Hash-Table 数据结构}
\label{lst:datastruct:Hash-Table Data Structures}
\end{listing}

\begin{figure}
\centering
\resizebox{3in}{!}{\includegraphics{datastruct/hashdiagram}}
\caption{Hash-Table 数据结构图}
\label{fig:datastruct:Hash-Table Data-Structure Diagram}
\end{figure}

\begin{fcvref}[ln:datastruct:hash_bkt:map_lock:map]
\Cref{lst:datastruct:Hash-Table Mapping and Locking}
展示了映射和加锁函数。
\Clnref{b,e}
展示了宏 \co{HASH2BKT()}，它将 hash 值映射到对应的 \co{ht_bucket} 结构。
该宏使用简单的取模运算：
如果需要更激进的 hash 处理，
调用者需要在从 key 映射到 hash 值时自行实现。
其余两个函数获取和释放与指定 hash 值对应的 \co{->htb_lock}。
\end{fcvref}

\begin{listing}
\input{CodeSamples/datastruct/hash/hash_bkt@map_lock.fcv}
\caption{Hash-Table 映射和加锁}
\label{lst:datastruct:Hash-Table Mapping and Locking}
\end{listing}

\begin{fcvref}[ln:datastruct:hash_bkt:lookup]
\Cref{lst:datastruct:Hash-Table Lookup}
展示了 \co{hashtab_lookup()}，
该函数返回指向具有指定 hash 和 key 的元素的指针，
如果该元素不存在则返回 \co{NULL}。
该函数同时接受 hash 值和指向 key 的指针，因为这允许该函数的用户使用任意 key 和
任意 hash 函数。
\Clnref{map} 将 hash 值映射到指向对应 hash bucket 的指针。
\clnrefrange{loop:b}{loop:e} 循环的每次迭代检查 bucket 的 hash chain 中的
一个元素。
\Clnref{hashmatch} 检查 hash 值是否匹配，如果不匹配，
\clnref{next}
继续检查下一个元素。
\Clnref{keymatch} 检查实际的 key 是否匹配，如果匹配，
\clnref{return} 返回指向匹配元素的指针。
如果没有元素匹配，\clnref{ret_NULL} 返回 \co{NULL}。
\end{fcvref}

\begin{listing}
\input{CodeSamples/datastruct/hash/hash_bkt@lookup.fcv}
\caption{Hash-Table 查找}
\label{lst:datastruct:Hash-Table Lookup}
\end{listing}

\QuickQuiz{
	\begin{fcvref}[ln:datastruct:hash_bkt:lookup]
	但在 key 能放入一个 unsigned long 的情况下，
	\cref{lst:datastruct:Hash-Table Lookup}中
	\clnrefrange{hashmatch}{return} 的双重比较不是效率低下吗？
	\end{fcvref}
}\QuickQuizAnswer{
	确实如此！
	然而，hash table 经常存储的信息的 key 是字符串之类的东西，
	不一定能放入一个 unsigned long 中。
	为 key 总能放入 unsigned long 的情况简化 hash table 实现，
	留作读者的练习。
}\QuickQuizEnd

\begin{listing}
\input{CodeSamples/datastruct/hash/hash_bkt@add_del.fcv}
\caption{Hash-Table 修改}
\label{lst:datastruct:Hash-Table Modification}
\end{listing}

\Cref{lst:datastruct:Hash-Table Modification}
展示了 \co{hashtab_add()} 和 \co{hashtab_del()} 函数，
分别用于向 hash table 添加和删除元素。

\begin{fcvref}[ln:datastruct:hash_bkt:add_del:add]
\co{hashtab_add()} 函数只是在 \clnref{set} 上设置元素的 hash 值，
然后在 \clnref{add:b,add:e} 上将其添加到对应的 bucket。
\end{fcvref}
\co{hashtab_del()} 函数只是将指定元素从其所在的 hash chain 中删除，
这得益于 hash chain 链表的双向链接特性。
在调用这两个函数之前，调用者需要确保没有其他线程正在访问
或修改同一个 bucket，例如，事先调用 \co{hashtab_lock()}。

\begin{listing}
\input{CodeSamples/datastruct/hash/hash_bkt@alloc_free.fcv}
\caption{Hash-Table 分配和释放}
\label{lst:datastruct:Hash-Table Allocation and Free}
\end{listing}

\Cref{lst:datastruct:Hash-Table Allocation and Free}
展示了 \co{hashtab_alloc()} 和 \co{hashtab_free()}，
分别执行 hash table 的分配和释放。
\begin{fcvref}[ln:datastruct:hash_bkt:alloc_free:alloc]
分配从 \clnrefrange{alloc:b}{alloc:e} 开始，分配底层内存。
如果 \clnref{chk_NULL} 检测到内存已耗尽，
\clnref{ret_NULL} 向调用者返回 \co{NULL}。
否则，\clnref{set_nbck,set_cmp} 初始化 bucket 数量和指向 key 比较函数的指针，
\clnrefrange{loop:b}{loop:e} 的循环初始化 bucket 本身，
包括 \clnref{init_head} 上的链表头和 \clnref{init_lock} 上的 lock。
最后，\clnref{return} 返回指向新分配的 hash table 的指针。
\end{fcvref}
\begin{fcvref}[ln:datastruct:hash_bkt:alloc_free:free]
\clnrefrange{b}{e} 上的 \co{hashtab_free()} 函数很直观。
\end{fcvref}

\subsection{哈希表性能}
\label{sec:datastruct:Hash-Table Performance}

\begin{figure}
\centering
\resizebox{2.5in}{!}{\includegraphics{CodeSamples/datastruct/hash/data/hps.perf.2020.11.26a/zoocpubktlin8}}
\caption{Schr\"odinger 动物园的只读 Hash-Table 性能}
\label{fig:datastruct:Read-Only Hash-Table Performance For Schroedinger's Zoo}
\end{figure}

在一个 2.1\,GHz Intel Xeon 系统的单个 28 核 socket 上，
使用具有 262,144 个 bucket 的 bucket 加锁 hash table 的性能结果如
\cref{fig:datastruct:Read-Only Hash-Table Performance For Schroedinger's Zoo}所示。
性能确实近似线性扩展，但远未达到理想水平，
即使只有 28 个 CPU。
这一差距的部分原因在于 lock 的获取和释放只在两个或更多 CPU 上才会产生
通信 cache miss。

\begin{figure}
\centering
\resizebox{2.5in}{!}{\includegraphics{CodeSamples/datastruct/hash/data/hps.perf.2020.11.26a/zoocpubktlin}}
\caption{Schr\"odinger 动物园的只读 Hash-Table 性能，448 个 CPU}
\label{fig:datastruct:Read-Only Hash-Table Performance For Schroedinger's Zoo; 448 CPUs}
\end{figure}

随着 CPU 数量增加，情况只会更糟，如
\cref{fig:datastruct:Read-Only Hash-Table Performance For Schroedinger's Zoo; 448 CPUs}所示。
我们甚至不需要展示理想性能：
29 个 CPU 及以上的性能令人惨不忍睹。
这清楚地说明了从少量 CPU 外推性能的危险性。

当然，性能崩溃的一个可能原因是需要更多的 hash bucket。
我们可以通过增加 hash bucket 数量来验证这一点。

\QuickQuiz{
	与其简单地增加 hash bucket 数量，
	对现有 hash bucket 进行 cache 对齐不是更好吗？
}\QuickQuizAnswer{
	答案取决于很多因素。
	如果 hash table 的每个 bucket 有大量元素，
	增加 hash bucket 数量显然更好。
	另一方面，如果 hash table 负载很轻，
	答案取决于硬件、hash 函数的有效性和工作负载。
	鼓励有兴趣的读者进行实验。
}\QuickQuizEnd


\begin{figure}
\centering
\resizebox{2.5in}{!}{\includegraphics{CodeSamples/datastruct/hash/data/hps.perf.2020.11.26a/zoocpubktsizelin}}
\caption{Schr\"odinger 动物园的只读 Hash-Table 性能，不同 Bucket 数量}
\label{fig:datastruct:Read-Only Hash-Table Performance For Schroedinger's Zoo; Varying Buckets}
\end{figure}

然而，如
\cref{fig:datastruct:Read-Only Hash-Table Performance For Schroedinger's Zoo; Varying Buckets}
所示，改变 bucket 数量几乎没有效果：
可扩展性仍然很差。
特别是，我们仍然看到在 29 个 CPU 及以上的急剧下降，
这清楚地证明了
\cpageref{sec:datastruct:NUMA-induced failure to scale}上复杂情况列表中
\cref{sec:datastruct:NUMA-induced failure to scale}提出的问题。
显然，还有其他原因。

问题在于这是一个多 socket 系统，CPU~0--27
和~224--251 映射到第一个 socket，如
\cref{fig:datastruct:NUMA Topology of System Under Test}所示。
因此，仅限于前 28 个 CPU 的测试运行表现相当好，
但涉及 socket~0 的 CPU~0--27 以及
socket~1 的 CPU~28 时，就会产生跨 socket 传输数据的开销。
正如\cref{sec:cpu:Hardware System Architecture}中所讨论的，
这会严重降低性能。
简而言之，大型多 socket 系统除了需要完全分区之外，还需要良好的引用局部性。
本章剩余部分将讨论在 hash table 本身内部提供良好引用局部性的方法，
但同时请注意，提供良好引用局部性的另一种方式是在 hash table 中放置大型数据元素。
例如，Schr\"odinger 可以通过在 hash table 的每个元素中放置动物的照片甚至视频来获得出色的 cache 局部性。
但对于需要包含小数据元素的 hash table 的人来说，请继续阅读！

\QuickQuiz{
	鉴于 Schr\"odinger 的动物园应用在跨 socket 时的负可扩展性，
	为什么不运行应用程序的多个副本，每个副本拥有动物的一个子集
	并限制在单个 socket 上运行呢？
}\QuickQuizAnswer{
	你完全可以这样做！
	事实上，你可以将这个想法扩展到大型集群系统，
	在集群的每个节点上运行应用程序的一个副本。
	这种做法称为 ``sharding''，在实践中被大型
	网络零售商大量使用~\cite{DeCandia:2007:DAH:1323293.1294281}。

	然而，如果你打算在多 socket 系统中按 socket 进行 sharding，
	为什么不购买单独的更小更便宜的单 socket 系统，
	然后在每个系统上运行数据库的一个分片呢？
}\QuickQuizEnd

这些 Schr\"odinger 动物园实验的一个关键特性是它们都是只读的。
这使得由 lock 获取引起的 cache miss 导致的性能下降更加令人痛苦。
即使我们没有更新底层的 hash table 本身，我们仍在为写入内存付出代价。
当然，如果 hash table 永远不会被更新，我们可以完全不用互斥。
这种方法非常直接，留给读者作为练习。
但即使偶尔有更新，避免写入也能避免 cache miss，
并允许读多数据在所有 cache 中复制，从而促进引用局部性。

因此，下一节将研究在更新罕见但可能随时发生的读多场景中可以进行的优化。

\setlength\dashlinedash{1pt}
\setlength\dashlinegap{2pt}

\begin{table}
\renewcommand*{\arraystretch}{1.2}
\centering
\begin{tabular}{r||c:c}
	& \multicolumn{2}{c}{Hyperthread} \\
	\cline{2-3}
	Socket & 0 &  1 \\
	\hline
	\hline
	0 &    0--27 & 224--251 \\
	\cdashline{2-3}
	1 &   28--55 & 252--279 \\
	\cdashline{2-3}
	2 &   56--83 & 280--307 \\
	\cdashline{2-3}
	3 &  84--111 & 308--335 \\
	\cdashline{2-3}
	4 & 112--139 & 336--363 \\
	\cdashline{2-3}
	5 & 140--167 & 364--391 \\
	\cdashline{2-3}
	6 & 168--195 & 392--419 \\
	\cdashline{2-3}
	7 & 196--223 & 420--447 \\
\end{tabular}
\caption{被测系统的 NUMA 拓扑}
\label{fig:datastruct:NUMA Topology of System Under Test}
\end{table}

\section{读多写少数据结构}
\label{sec:datastruct:Read-Mostly Data Structures}
%
\epigraph{对症下药。}{中国谚语}

虽然分区数据结构可以提供优秀的可扩展性，
但 \IXacr{numa} 效应可能导致性能和可扩展性的严重下降。
此外，读侧同步的需要可能在以读为主的场景下降低性能。
然而，我们可以使用 RCU 来同时实现性能和可扩展性，
RCU 已在\cref{sec:defer:Read-Copy Update (RCU)}中介绍。
使用 \IXpl{hazard pointer}（\path{hazptr.c}）~\cite{MagedMichael04a}
也可以实现类似的结果，本节展示的性能结果中将包含
hazard pointer 的数据~\cite{McKenney:2013:SDS:2483852.2483867}。

\subsection{RCU 保护的哈希表实现}
\label{sec:datastruct:RCU-Protected Hash Table Implementation}

对于使用每 bucket 加锁的 RCU 保护 hash table，
更新者使用如\cref{sec:datastruct:Partitionable Data Structures}所示的加锁方式，
而读取者使用 RCU\@。
数据结构保持如\cref{lst:datastruct:Hash-Table Data Structures}所示，
\co{HASH2BKT()}、\co{hashtab_lock()} 和 \co{hashtab_unlock()}
函数保持如\cref{lst:datastruct:Hash-Table Mapping and Locking}所示。
然而，读取者使用更轻量级的并发控制，
即\cref{lst:datastruct:RCU-Protected Hash-Table Read-Side Concurrency Control}
中所示的 \co{hashtab_lock_lookup()} 和 \co{hashtab_unlock_lookup()}。

\begin{listing}
\input{CodeSamples/datastruct/hash/hash_bkt_rcu@lock_unlock.fcv}
\caption{RCU 保护的 Hash-Table 读侧并发控制}
\label{lst:datastruct:RCU-Protected Hash-Table Read-Side Concurrency Control}
\end{listing}

\Cref{lst:datastruct:RCU-Protected Hash-Table Lookup}
展示了 RCU 保护的每 bucket 加锁 hash table 的 \co{hashtab_lookup()}。
它与\cref{lst:datastruct:Hash-Table Lookup}中的版本完全相同，
只是将 \apiur{cds_list_for_each_entry()} 替换为了
\apiur{cds_list_for_each_entry_rcu()}。
这两个原语都遍历由 \co{htb->htb_head} 引用的 hash chain，
但 \co{cds_list_for_each_entry_rcu()} 还能在并发插入的情况下正确地保证内存顺序。
这是两种 hash table 实现之间的一个重要区别：
与纯粹的每 bucket 加锁实现不同，RCU 保护的实现允许查找与插入和删除并发运行，
因此需要像 \co{cds_list_for_each_entry_rcu()} 这样的 RCU 感知原语
来正确处理这种额外的并发。
还要注意，\co{hashtab_lookup()} 的调用者必须在 RCU 读侧临界区内，
例如，调用者必须在调用 \co{hashtab_lookup()} 之前调用
\co{hashtab_lock_lookup()}（当然之后还要调用 \co{hashtab_unlock_lookup()}）。

\begin{listing}
\input{CodeSamples/datastruct/hash/hash_bkt_rcu@lookup.fcv}
\caption{RCU 保护的 Hash-Table 查找}
\label{lst:datastruct:RCU-Protected Hash-Table Lookup}
\end{listing}

\QuickQuiz{
	但如果 hash table 中的元素可以与查找并发删除，
	这不意味着查找可能返回一个刚被删除的数据元素的引用吗？
}\QuickQuizAnswer{
	确实如此！
	这就是为什么 \co{hashtab_lookup()} 必须在 RCU 读侧临界区内调用，
	也是为什么 \co{hashtab_add()} 和 \co{hashtab_del()} 也必须使用
	RCU 感知的链表操作原语。
	最后，这也是为什么 \co{hashtab_del()} 的调用者必须等待一个
	grace period（例如，通过调用 \co{synchronize_rcu()}）
	才能释放被删除的元素。
	这将确保所有可能引用新删除元素的 RCU 读取者
	在该元素被释放之前已经完成。
}\QuickQuizEnd

\begin{listing}
\input{CodeSamples/datastruct/hash/hash_bkt_rcu@add_del.fcv}
\caption{RCU 保护的 Hash-Table 修改}
\label{lst:datastruct:RCU-Protected Hash-Table Modification}
\end{listing}

\Cref{lst:datastruct:RCU-Protected Hash-Table Modification}
展示了 \co{hashtab_add()} 和 \co{hashtab_del()}，
两者都与\cref{lst:datastruct:Hash-Table Modification}中
非 RCU hash table 的对应函数非常相似。
\co{hashtab_add()} 函数使用 \apiur{cds_list_add_rcu()} 而不是
\apiur{cds_list_add()}，以确保在向 hash table 添加元素的同时
有人查找时的正确排序。
\co{hashtab_del()} 函数使用 \apiur{cds_list_del_rcu()} 而不是
\apiur{cds_list_del_init()}，以处理元素在被删除之前刚好被查找的情况。
与 \co{cds_list_del_init()} 不同，\co{cds_list_del_rcu()} 保留前向指针不变，
这样 \co{hashtab_lookup()} 就可以遍历到新删除元素的后继者。

当然，在调用 \co{hashtab_del()} 之后，调用者必须等待一个
RCU \IX{grace period}（例如，通过调用 \co{synchronize_rcu()}），
然后才能释放或重用新删除元素的内存。

\subsection{RCU 保护的哈希表验证}
\label{sec:datastruct:RCU-Protected Hash Table Validation}

虽然验证这一主题在\cref{chp:Validation}中有详细介绍，
但事实上，具有无锁（lock-free） RCU 保护查找的 hash table 需要尽早对验证给予特别关注。

测试套件（``\path{hashtorture.h}''）包含一个 \co{smoketest()} 函数，
用于验证一系列特定的单线程添加、删除和查找操作是否给出预期结果。

并发测试运行让每个更新者线程控制其负责的那部分元素，
这允许 assertion 检查以下问题：

\begin{enumerate}
\item	一个即将添加的元素根据 \co{hastab_lookup()} 已经存在于表中。
\item	一个即将添加的元素的 \co{->in_table} 标志显示其已在表中。
\item	一个即将删除的元素根据 \co{hastab_lookup()} 不在表中。
\item	一个即将删除的元素的 \co{->in_table} 标志显示其不在表中。
\end{enumerate}

此外，并发测试运行将查找与更新并发执行，以捕获各种数据结构损坏问题。
一些测试运行还会在查找和更新的同时持续调整 hash table 大小，
以验证正确的行为，并验证调整大小不会过度延迟读取者或更新者。

最后，并发测试输出的统计信息可用于追踪性能和可扩展性问题，
这提供了\cref{sec:datastruct:RCU-Protected Hash Table Performance}所用的原始数据。

\QuickQuiz{
	\path{hashtorture.h} 文件包含超过 1,000 行！
	这算是一个全面的测试吗？
}\QuickQuizAnswer{
	不算。

	\path{hashtorture.h} 的测试是一个好的开始，对于教科书算法来说足够了。
	如果这段代码要用于生产环境，则需要更多的测试：

	\begin{enumerate}
	\item	有一部分元素始终驻留在表中，并验证无论有多少以及何种类型的
		并发更新正在进行，查找都能找到这些元素。
	\item	将一个更新者与一个或多个读取者配对，验证在添加一个元素后，
		一旦读取者成功查找到该元素，后续所有查找都成功。
		"后续"的定义取决于表的一致性要求。
	\item	将一个更新者与一个或多个读取者配对，验证在删除一个元素后，
		一旦读取者对该元素的查找失败，后续所有查找也都失败。
	\end{enumerate}

	还有很多类似的测试，具体性质取决于你特定 hash table 的需求细节。
}\QuickQuizEnd

所有代码都需要大量的验证工作，而高性能并发代码比大多数代码需要更多的验证。

\subsection{RCU 保护的哈希表性能}
\label{sec:datastruct:RCU-Protected Hash Table Performance}

\begin{figure}
\centering
\resizebox{2.5in}{!}{\includegraphics{CodeSamples/datastruct/hash/data/hps.perf.2020.11.26a/zoocpu}}
\caption{Schr\"odinger 动物园的只读 RCU 保护 Hash-Table 性能}
\label{fig:datastruct:Read-Only RCU-Protected Hash-Table Performance For Schroedinger's Zoo}
\end{figure}

\Cref{fig:datastruct:Read-Only RCU-Protected Hash-Table Performance For Schroedinger's Zoo}
展示了 RCU 保护和 hazard pointer 保护的 hash table 与上一节的每 bucket 加锁实现的
只读性能对比。
如你所见，RCU 和 hazard pointer 的性能和可扩展性都远优于每 bucket 加锁，
因为只读复制避免了 NUMA 效应。
差异随线程数量增加而增大。
还展示了全局加锁实现的结果，不出所料，
其结果甚至比每 bucket 加锁实现更差。
RCU 比 hazard pointer 略好。

\begin{figure}
\centering
\resizebox{2.5in}{!}{\includegraphics{CodeSamples/datastruct/hash/data/hps.perf.2020.11.26a/zoocpulin}}
\caption{Schr\"odinger 动物园的只读 RCU 保护 Hash-Table 性能，线性刻度}
\label{fig:datastruct:Read-Only RCU-Protected Hash-Table Performance For Schroedinger's Zoo; Linear Scale}
\end{figure}

\Cref{fig:datastruct:Read-Only RCU-Protected Hash-Table Performance For Schroedinger's Zoo; Linear Scale}
以线性刻度展示了相同的数据。
这使全局加锁的曲线落入 x 轴，但能更容易辨别
RCU 和 hazard pointer 的非理想性能。
两者在 224 个 CPU 处都显示了斜率变化，这是由硬件多线程引起的。
在 224 个及更少的 CPU 时，每个线程独占一个核心。
在这种情况下，RCU 比 hazard pointer 表现更好，因为后者的读侧
\IXpl{memory barrier}导致核心内的空闲时间。
简而言之，从单个硬件线程来看，RCU 比 hazard pointer 能更好地利用一个核心。

这种情况在 224 个 CPU 以上发生了变化。
因为 RCU 从单个硬件线程就使用了每个核心一半以上的资源，
RCU 从每个核心的第二个硬件线程获得的收益相对较少。
Hazard pointer 的曲线斜率在 224 个 CPU 处也下降了，但不那么明显，
因为第二个硬件线程能够填补第一个硬件线程因 \IXh{memory-barrier}{latency}
而停滞的时间。
我们将在后面的章节中看到，第二个硬件线程的这一优势取决于工作负载。

\begin{figure}
\centering
\resizebox{2.5in}{!}{\includegraphics{CodeSamples/datastruct/hash/data/hps.perf.2020.11.26a/zoocpulinqsbr}}
\caption{Schr\"odinger 动物园的只读 RCU 保护 Hash-Table 性能（含 QSBR），线性刻度}
\label{fig:datastruct:Read-Only RCU-Protected Hash-Table Performance For Schroedinger's Zoo including QSBR; Linear Scale}
\end{figure}

但为什么 RCU 的性能比理想值低五倍？
一个可能是 \co{rcu_read_lock()} 和 \co{rcu_read_unlock()} 操作的每线程计数器
拖慢了速度。
因此\cref{fig:datastruct:Read-Only RCU-Protected Hash-Table Performance For Schroedinger's Zoo including QSBR; Linear Scale}
添加了 RCU 的 \IXacr{qsbr} 变体的结果，其读侧原语什么都不做。
虽然 QSBR 确实比 RCU 表现稍好，但仍然比理想值低约五倍。

\begin{figure}
\centering
\resizebox{2.5in}{!}{\includegraphics{CodeSamples/datastruct/hash/data/hps.perf.2020.11.26a/zoocpulinqsbrunsync}}
\caption{Schr\"odinger 动物园的只读 RCU 保护 Hash-Table 性能（含 QSBR 和无同步），线性刻度}
\label{fig:datastruct:Read-Only RCU-Protected Hash-Table Performance For Schroedinger's Zoo including QSBR and Unsynchronized; Linear Scale}
\end{figure}

\Cref{fig:datastruct:Read-Only RCU-Protected Hash-Table Performance For Schroedinger's Zoo including QSBR and Unsynchronized; Linear Scale}
添加了完全无同步的结果，这是可行的，因为这是一个只读基准测试，没有需要同步的内容。
即使完全没有任何同步，性能仍然远未达到理想水平，从而证明了
\cpageref{sec:datastruct:Cache-size-induced failure to scale redux}上
复杂情况列表中\cref{sec:datastruct:Cache-size-induced failure to scale,sec:datastruct:Cache-size-induced failure to scale redux}的问题。

问题在于这个系统的 socket 有 28 个核心，其 cache 大小相对较小，如
\cpageref{tab:cpu:Cache Geometry for 8-Socket System With Intel Xeon Platinum 8176 CPUs @ 2.10GHz}上的
\cref{tab:cpu:Cache Geometry for 8-Socket System With Intel Xeon Platinum 8176 CPUs @ 2.10GHz}所示。
每个 hash bucket（\co{struct ht_bucket}）占用 56 字节，
在 RCU 和 QSBR 运行中每个元素（\co{struct zoo_he}）占用 72 字节。
生成\cref{fig:datastruct:Read-Only RCU-Protected Hash-Table Performance For Schroedinger's Zoo including QSBR and Unsynchronized; Linear Scale}
的基准测试使用了 262,144 个 bucket 和最多 262,144 个元素，
总计 33,554,448 字节，这不仅超出 1,048,576 字节的 L2 cache 三十多倍，
而且也令人不安地接近 40,370,176 字节的 L3 cache 大小，
尤其考虑到该 cache 只有 11 路关联。
这意味着 L2 cache 冲突将成为常态，L3 cache 冲突也不少见，
由此产生的 cache miss 将降低性能。
在这种情况下，瓶颈不在 CPU，而在内存系统。
% This data was generated wtih 262,144 buckets and elements.
% RCU hash ht_elem structure is 40 bytes plus 32 more from zoo_he for 72.
% struct ht_bucket is 56 bytes.
% For "# B" RCU runs, hash buckets weigh in at 14,680,080 bytes and the hash
%    elements up to 18,874,368 bytes.  The total is 33,554,448 bytes, which
%    in theory fits into the 40,370,176-byte L3 cache, but definitely not
%    the 1,048,576-byte L2 cache, much less the 32,768-byte L0 and L1 caches.
%    The L3 is 11-way set associative.
%    So we have 8192 4K pages of data (buckets and elements) that must
%    fit into 896 11-way set-associated cache sets.

通过检查无同步代码可以找到这种内存系统瓶颈的更多证据。
该代码不需要 lock，因此每个 hash bucket 只占 16 字节，
而 RCU 和 QSBR 需要 56 字节。
类似地，每个 hash table 元素只占 56 字节，
而 RCU 和 QSBR 需要 72 字节。
因此，单 CPU 无同步运行的性能比 QSBR 或 RCU 快约一半并不令人惊讶。

\QuickQuiz{
	我们怎么能如此确定问题出在 hash table 大小上，
	尤其是考虑到
	\cpageref{fig:datastruct:Read-Only Hash-Table Performance For Schroedinger's Zoo; Varying Buckets}上的
	\cref{fig:datastruct:Read-Only Hash-Table Performance For Schroedinger's Zoo; Varying Buckets}
	显示改变 hash table 大小几乎没有效果？
	问题会不会是 \IX{false sharing} 之类的？
}\QuickQuizAnswer{
	好问题！

	False sharing 需要写操作，而在这个只查找基准测试的无同步和 RCU 运行中
	没有写操作。
	因此问题不是 false sharing。

\begin{figure}
\centering
\resizebox{3in}{!}{\includegraphics{CodeSamples/datastruct/hash/data/hps.perf-hashsize.2020.12.29a/zoohashsize}}
\caption{Schr\"odinger 动物园在 448 个 CPU 下的只读 RCU 保护 Hash-Table 性能，不同表大小}
\label{fig:datastruct:Read-Only RCU-Protected Hash-Table Performance For Schr\"odinger's Zoo at 448 CPUs; Varying Table Size}
\end{figure}

	还不信？
	那就看看
	\cref{fig:datastruct:Read-Only RCU-Protected Hash-Table Performance For Schr\"odinger's Zoo at 448 CPUs; Varying Table Size}
	中的双对数图，它展示了 448 个 CPU 下性能随 hash table 大小
	（即 bucket 数量和最大元素数量）的变化。
	大小为 1,024 的 hash table 有 1,024 个 bucket，
	最多包含 1,024 个元素，平均占用为 512 个元素。
	因为这是一个只读基准测试，实际占用始终等于平均占用。

	该图显示在约 8,000 个元素以下时性能接近理想，
	即当 hash table 的数据量不到 1\,MB 时。
	这种接近理想的性能与
	\cpageref{fig:defer:Pre-BSD Routing Table Protected by RCU}上的
	\cref{fig:defer:Pre-BSD Routing Table Protected by RCU}
	所示的 pre-BSD 路由表一致，即使在 448 个 CPU 下也是如此。
	然而，在约 8,000 个元素处性能显著下降（这是双对数图），
	这正是 1,048,576 字节 L2 cache 溢出的地方。
	在约 300,000 个元素处性能急剧下降（即使在双对数图上也是如此），
	这是 40,370,176 字节 L3 cache 溢出的地方。
	这表明内存系统瓶颈是深刻的，
	对于大型 hash table，性能下降远超一个数量级。
	这不应该令人惊讶，因为大小为 8,388,608 的 hash table
	占用约 1\,GB 内存，超出 L3 cache 25 倍。

	\cref{fig:datastruct:Read-Only Hash-Table Performance For Schroedinger's Zoo; Varying Buckets}
	（\cpageref{fig:datastruct:Read-Only Hash-Table Performance For Schroedinger's Zoo; Varying Buckets}）
	显示影响很小的原因是其数据来自 bucket 加锁的 hash table，
	其中加锁开销和竞争掩盖了 cache 容量效应。
	相比之下，RCU 和 hazard pointer 的读取者避免了对共享数据的存储，
	这意味着 cache 容量效应变得突出。

	还不满意？
	找一个多 socket 系统来运行这段代码，利用可用的性能计数器硬件。
	这些硬件应该能让你追踪在你特定系统上出现的任何性能下降的精确原因。
	通过做这个练习获得的经验将非常宝贵，
	使你比那些仅从理论上理解这个问题的人拥有显著优势。\footnote{
		当然，理论理解比完全不理解要好。
		至少大多数时候是这样。}
}\QuickQuizEnd

如果进一步减小内存占用呢？
\cpageref{fig:defer:Pre-BSD Routing Table Protected by RCU QSBR With Non-Initial rcu-head}上的
\cref{fig:defer:Pre-BSD Routing Table Protected by RCU QSBR With Non-Initial rcu-head}
显示 RCU 在 pre-BSD 路由表所代表的小得多的数据结构上
达到了非常接近理想的性能。

\QuickQuiz{
	内存系统在这个大型系统上是一个严重的瓶颈。
	为什么要在一个没有给它们足够内存带宽来做有用工作的系统上
	安装 448 个 CPU？
}\QuickQuizAnswer{
	将这个大型昂贵的系统用于仅由小数据元素的简单 hash table 查找组成的
	工作负载确实不是个好主意。
	然而，这个系统对于许多具有更多处理和更少内存访问的工作负载极其有用。
	例如，一些内存数据库在这类系统上运行得非常好，
	尽管运行的查询集比本章中的基准测试复杂得多。
	例如，这类系统可能在处理存储在每个元素中的图像或视频流，
	由于由此产生的顺序内存访问比纯粹的指针追踪工作负载
	能更好地利用可用的内存带宽，从而提供额外的性能优势。

	但让这成为一个教训。
	现代计算机系统有各种各样的形状和大小，
	选择一个适合你应用的系统经常需要非常谨慎。
	而且也许更常见的是，需要大量的精力和工作
	来调整你的应用以适应手头的特定计算机系统。
}\QuickQuizEnd

\begin{figure}
\centering
\resizebox{2.5in}{!}{\includegraphics{CodeSamples/datastruct/hash/data/hps.perf.2020.11.26a/zoocatonly}}
\caption{Schr\"odinger 动物园在 64 个 CPU 下的仅查找猫的读侧 RCU 保护 Hash-Table 性能}
\label{fig:datastruct:Read-Side Cat-Only RCU-Protected Hash-Table Performance For Schroedinger's Zoo at 64 CPUs}
\end{figure}

如前所述，Schr\"odinger 对他的猫的受欢迎程度感到
惊讶~\cite{ErwinSchroedinger1935Cat}，
但认识到需要在他的设计中反映这种受欢迎程度。
\Cref{fig:datastruct:Read-Side Cat-Only RCU-Protected Hash-Table Performance For Schroedinger's Zoo at 64 CPUs}
展示了 64 CPU 运行的结果，变化了只做查找猫操作的 CPU 数量。
RCU 和 hazard pointer 都能很好地应对这一挑战，
但 bucket 加锁的可扩展性为负，最终表现得和全局加锁一样差。
这不应令人惊讶，因为如果所有 CPU 只做查找猫的操作，
猫所在 bucket 对应的 lock 实际上就是一个全局 lock。

这个仅查找猫的基准测试说明了完全分区 sharding 方法的一个潜在问题。
只有与猫的分区关联的 CPU 才能访问猫，限制了仅查找猫的吞吐量。
当然，很多应用具有良好的负载分散特性，
对于这些应用 sharding 工作得很好。
然而，sharding 不能很好地处理 ``\IXpl{hot spot}''，
Schr\"odinger 的猫所代表的热点就是一个典型例子。

\begin{figure}
\centering
\resizebox{2.5in}{!}{\includegraphics{CodeSamples/datastruct/hash/data/hps.perf.2020.11.26a/zooupdatelu}}
\caption{有更新时 Schr\"odinger 动物园的读侧 RCU 保护 Hash-Table 性能}
\label{fig:datastruct:Read-Side RCU-Protected Hash-Table Performance For Schroedinger's Zoo in the Presence of Updates}
\end{figure}

如果我们只打算读取数据，那么一开始就不需要任何并发控制。
\Cref{fig:datastruct:Read-Side RCU-Protected Hash-Table Performance For Schroedinger's Zoo in the Presence of Updates}
因此展示了更新对读取者的影响。
在该图的最左端，除一个 CPU 外所有 CPU 都在执行查找，
而在右端所有 448 个 CPU 都在执行更新。
对于所有四种实现，每毫秒查找次数随更新 CPU 数量增加而减少，
当然在所有 448 个 CPU 都在更新时达到每毫秒零次查找。
Hazard pointer 和 RCU 比每 bucket 加锁表现好，
因为它们的读取者不会增加更新侧的 lock 竞争。
随着更新者数量增加，RCU 相对于 hazard pointer 表现更好，
这是因为后者的读侧 memory barrier 产生了更大的开销，
尤其是在有更新存在时，特别是当执行涉及多个 socket 时。
因此，现代硬件很可能对 memory barrier 的执行进行了大量优化，
在只读情况下大幅降低了 memory barrier 开销。
% @@@ When there is a section covering trees, reference sharding with
% @@@ CPUs being permitted to directly look up popular entries.

\begin{figure}
\centering
\resizebox{2.5in}{!}{\includegraphics{CodeSamples/datastruct/hash/data/hps.perf.2020.11.26a/zooupdate}}
\caption{Schr\"odinger 动物园的更新侧 RCU 保护 Hash-Table 性能}
\label{fig:datastruct:Update-Side RCU-Protected Hash-Table Performance For Schroedinger's Zoo}
\end{figure}

\cref{fig:datastruct:Read-Side RCU-Protected Hash-Table Performance For Schroedinger's Zoo in the Presence of Updates}
展示了增加更新速率对查找的影响，而
\cref{fig:datastruct:Update-Side RCU-Protected Hash-Table Performance For Schroedinger's Zoo}
展示了增加更新速率对更新本身的影响。
同样，在图的左端除一个 CPU 外所有 CPU 都在执行查找，
在图的右端所有 448 个 CPU 都在执行更新。
Hazard pointer 和 RCU 一开始具有显著优势，因为与 bucket 加锁不同，
读取者不会排斥更新者。
然而，随着更新 CPU 数量增加，更新侧的延迟执行开销开始显现，
首先是 RCU，然后是 hazard pointer。
当然，这三种实现都胜过全局加锁，而且差距超过一个数量级。

在\cref{fig:datastruct:Read-Side RCU-Protected Hash-Table Performance For Schroedinger's Zoo in the Presence of Updates}
中观察到的查找性能差异很可能受到更新速率差异的影响。
检验这一点的一种方法是人为限制每 bucket 加锁和 hazard pointer 的更新速率
以匹配 RCU 的更新速率。
这样做并不能显著改善每 bucket 加锁的查找性能，
也不能缩小 hazard pointer 和 RCU 之间的差距。
然而，从 hazard pointer 中移除读侧 memory barrier
（从而产生一个不安全的实现）确实几乎消除了
hazard pointer 和 RCU 之间的差距。
虽然这种不安全的 hazard pointer 实现通常对基准测试来说够可靠，
但绝对不建议在生产环境中使用。

\QuickQuiz{
	从 28 个 CPU 外推到 448 个 CPU 的危险性在
	\cref{sec:datastruct:Hash-Table Performance}中已经很清楚了。
	从 448 个 CPU 向上外推会更安全吗？
}\QuickQuizAnswer{
	理论上不会更安全，在更大的系统上运行这些程序将是一个有用的练习。
	实际上，拥有超过 448 个 CPU 的系统非常少，
	而拥有超过 28 个 CPU 的系统则数量巨大。
	这意味着虽然外推超过 448 个 CPU 是危险的，但很少需要这样做。

	此外，其他测试表明 RCU 读侧原语在至少 1024 个 CPU 的情况下
	提供一致的性能和可扩展性。
	然而，回顾
	\cref{fig:datastruct:Read-Only RCU-Protected Hash-Table Performance For Schr\"odinger's Zoo at 448 CPUs; Varying Table Size}
	及其相关评论是有用的。
	要知道，与提供这些数据的 448 CPU 系统不同，
	享有线性可扩展性直到 1024 个 CPU 的系统拥有出色的内存带宽。
}\QuickQuizEnd

这种情况证明了
\cpageref{sec:datastruct:Update-activity failure to scale}上
复杂情况列表中\cref{sec:datastruct:Update-activity failure to scale}的问题。

% @@@ Testing strategy.  Summarize hashtorture, add QQ for additional
% @@@ things that could be done.  Note difference between textbook
% @@@ and production algorithms, also the need for clear understanding
% @@@ of what level of consistency is required.

\subsection{RCU 保护的哈希表讨论}
\label{sec:datastruct:RCU-Protected Hash Table Discussion}

RCU 和 hazard pointer 实现的一个后果是，
一对并发读取者可能对猫的状态产生分歧。
例如，其中一个读取者可能在猫的数据结构被移除之前获取了指向它的指针，
而另一个读取者可能在之后获取了这个指针。
第一个读取者会认为猫是活的，而第二个读取者会认为猫已经死了。

这种情况完全适合 Schr\"odinger 的猫，但事实证明，
对于普通的非量子猫也是完全合理的。
毕竟，精确确定一只动物何时出生或死亡是不可能的。

为了理解这一点，假设我们通过心跳来检测猫的死亡。
这就引出了一个问题：在最后一次心跳之后我们应该等多久才能宣布死亡。
只等一毫秒显然是荒谬的，因为那样一只健康的活猫将不得不
被宣布死亡——然后复活——每秒超过一次。
等整整一个月同样荒谬，因为到那时可怜的猫的死亡
对嗅觉传感器来说将毫无疑问地明显。

\begin{figure}
\centering
\resizebox{3in}{!}{\includegraphics{cartoons/2013-08-is-it-dead}}
\caption{连兽医也意见不一！}
\ContributedBy{fig:datastruct:Even Veterinarians Disagree}{Melissa Broussard}
\end{figure}

因为动物的心脏可能停跳几秒钟然后再次跳动，
所以在及时识别死亡和误报概率之间存在权衡。
一对兽医完全可能对最后一次心跳到宣布死亡之间应等待的时间意见不一。
例如，一位兽医可能在最后一次心跳后三十秒宣布死亡，
而另一位可能坚持等整整一分钟。
在这种情况下，两位兽医将在最后一次心跳后的第二个三十秒内
对猫的状态意见不一，
如\cref{fig:datastruct:Even Veterinarians Disagree}所富有想象力地描绘的那样。

\pplsur{Weiner}{Heisenberg} 教会了我们与这种不确定性共处~\cite{WeinerHeisenberg1927Uncertain}，
这是件好事，因为计算硬件和软件的行为也类似。
例如，你怎么知道一个计算硬件已经故障了？
通常是因为它没有及时响应。
就像猫的心跳一样，这导致了一个不确定性窗口，
不知道硬件是真的故障了还是只是慢了。

此外，大多数计算系统旨在与外部世界交互。
因此，与外部世界的一致性至关重要。
然而，如我们在
\cpageref{fig:defer:Response Time of RCU vs. Reader-Writer Locking}上的
\cref{fig:defer:Response Time of RCU vs. Reader-Writer Locking}中看到的，
增强内部一致性可能以降低外部一致性为代价。
RCU 和 hazard pointer 等技术放弃了一定程度的内部一致性
以获得更好的外部一致性。

简而言之，内部一致性不一定是所有问题领域的自然组成部分，
而且通常在性能、可扩展性、与外部世界的一致性
~\cite{AndreasHaas2012FIFOisnt,AndreasHaas2013CFRelaxedQueues,10.5555/3241639.3241645}
或以上所有方面付出巨大代价。

\section{不可分区数据结构}
\label{sec:datastruct:Non-Partitionable Data Structures}
%
\epigraph{如果需要迈出一大步，不要害怕。
	  你不能用两小步跨过深渊。}
	 {David Lloyd George}

固定大小的 hash table 是完全可分区的，但可调整大小的 hash table 在增长或缩小时
会带来分区挑战，如\cref{fig:datastruct:Partitioning Problems}所富有想象力地描绘的那样。
然而，事实证明可以构建高性能可扩展的 RCU 保护 hash table，
如以下各节所述。

\begin{figure}
\centering
\resizebox{3in}{!}{\includegraphics{cartoons/2014_Hash-table-hydra}}
\caption{分区问题}
\ContributedBy{fig:datastruct:Partitioning Problems}{Melissa Broussard}
\end{figure}

\subsection{可调整大小的哈希表设计}
\label{sec:datastruct:Resizable Hash Table Design}

与 2000 年代初的情况形成令人愉快的对比，现在有几种不同类型的
可扩展 RCU 保护 hash table。
第一种（也是最简单的）是由 \ppl{Herbert}{Xu}
为 Linux 内核开发的~\cite{HerbertXu2010RCUResizeHash}，
将在以下各节中描述。
另外两种在\cref{sec:datastruct:Other Resizable Hash Tables}中简要介绍。

第一种 hash table 实现背后的关键洞察是，
每个数据元素可以有两组链表指针，其中一组当前正被 RCU 读取者使用
（以及非 RCU 更新者），另一组用于构建新的调整大小后的 hash table。
这种方法允许查找、插入和删除操作都与调整大小操作（以及彼此）并发运行。

\begin{figure}
\centering
\resizebox{3in}{!}{\includegraphics{datastruct/hashxu-a}}
\caption{增长双链表 Hash Table，状态~(a)}
\label{fig:datastruct:Growing a Two-List Hash Table; State (a)}
\end{figure}

\begin{figure}
\centering
\resizebox{3in}{!}{\includegraphics{datastruct/hashxu-b}}
\caption{增长双链表 Hash Table，状态~(b)}
\label{fig:datastruct:Growing a Two-List Hash Table; State (b)}
\end{figure}

调整大小操作的过程如
\crefrange{fig:datastruct:Growing a Two-List Hash Table; State (a)}
{fig:datastruct:Growing a Two-List Hash Table; State (d)}所示，
初始的两 bucket 状态如
\cref{fig:datastruct:Growing a Two-List Hash Table; State (a)}所示，
时间从一个图前进到下一个图。
初始状态使用零索引链接将元素串入 hash bucket。
分配一个四 bucket 数组，使用一索引链接将元素串入这四个新的 hash bucket。
这导致了\cref{fig:datastruct:Growing a Two-List Hash Table; State (b)}
所示的状态~(b)，读取者仍在使用原来的两 bucket 数组。

\begin{figure}
\centering
\resizebox{3in}{!}{\includegraphics{datastruct/hashxu-c}}
\caption{增长双链表 Hash Table，状态~(c)}
\label{fig:datastruct:Growing a Two-List Hash Table; State (c)}
\end{figure}

\begin{figure}
\centering
\resizebox{3in}{!}{\includegraphics{datastruct/hashxu-d}}
\caption{增长双链表 Hash Table，状态~(d)}
\label{fig:datastruct:Growing a Two-List Hash Table; State (d)}
\end{figure}

新的四 bucket 数组暴露给读取者，然后 grace period 操作等待所有读取者完成，
产生状态~(c)，如\cref{fig:datastruct:Growing a Two-List Hash Table; State (c)}所示。
在这个状态下，所有读取者都在使用新的四 bucket 数组，
这意味着旧的两 bucket 数组现在可以被释放，产生状态~(d)，
如\cref{fig:datastruct:Growing a Two-List Hash Table; State (d)}所示。

这一设计带来了相对直接的实现，这是下一节的主题。

\subsection{可调整大小的哈希表实现}
\label{sec:datastruct:Resizable Hash Table Implementation}

\begin{fcvref}[ln:datastruct:hash_resize:data]
调整大小是通过经典的插入间接层方法完成的，
在这种情况下是 \co{ht} 结构，如
\cref{lst:datastruct:Resizable Hash-Table Data Structures}
（\path{hash_resize.c}）的
\clnrefrange{ht:b}{ht:e} 所示。
\clnrefrange{hashtab:b}{hashtab:e} 所示的 \co{hashtab} 结构仅包含
一个指向当前 \co{ht} 结构的指针以及一个用于序列化并发调整大小尝试的 spinlock。
如果我们使用传统的基于 lock 或原子操作（atomic operation）的实现，
这个 \co{hashtab} 结构可能会从性能和可扩展性两方面成为严重的瓶颈。
然而，由于调整大小操作应该相对不频繁，
我们应该能够很好地利用 RCU\@。

\begin{listing}
\input{CodeSamples/datastruct/hash/hash_resize@data.fcv}
\caption{可调整大小的 Hash-Table 数据结构}
\label{lst:datastruct:Resizable Hash-Table Data Structures}
\end{listing}

\co{ht} 结构表示 hash table 的特定大小，
由 \clnref{ht:nbuckets} 上的 \co{->ht_nbuckets} 字段指定。
大小存储在包含 bucket 数组（\clnref{ht:bkt} 上的 \co{->ht_bkt[]}）的
同一结构中，以避免大小和数组之间的不匹配。
\clnref{ht:resize_cur} 上的 \co{->ht_resize_cur} 字段等于 $-1$，
除非正在进行调整大小操作，
在这种情况下它表示其元素正在被插入新 hash table 的 bucket 索引，
新 hash table 由 \clnref{ht:new} 上的 \co{->ht_new} 字段引用。
如果没有正在进行的调整大小操作，\co{->ht_new} 为 \co{NULL}。
因此，调整大小操作通过分配新的 \co{ht} 结构并通过 \co{->ht_new} 指针引用它，
然后推进 \co{->ht_resize_cur} 遍历旧表的 bucket 来进行。
当所有元素都添加到新表后，新表被链接到 \co{hashtab} 结构的 \co{->ht_cur} 字段。
一旦所有旧读取者完成，旧 hash table 的 \co{ht} 结构就可以被释放。

\clnref{ht:idx} 上的 \co{->ht_idx} 字段表示 hash table 的这个实例
正在使用两组链表指针中的哪一组，
用于索引 \co{ht_elem} 结构中 \clnref{ht_elem:next} 上的
\co{->hte_next[]} 数组。

\clnrefrange{ht:cmp}{ht:getkey} 上的
\co{->ht_cmp()}、\co{->ht_gethash()} 和 \co{->ht_getkey()} 字段
共同定义了每个元素的 key 和 hash 函数。
\co{->ht_cmp()} 函数将指定的 key 与指定元素的 key 进行比较，
\co{->ht_gethash()} 计算指定 key 的 hash，
\co{->ht_getkey()} 从包含的数据元素中提取 key。

\clnrefrange{hls:b}{hls:e} 所示的 \co{ht_lock_state}
用于将 lock 状态从新的 \co{hashtab_lock_mod()}
传递到 \co{hashtab_add()}、\co{hashtab_del()} 和 \co{hashtab_unlock_mod()}。
该状态防止算法在并发调整大小操作期间被重定向到错误的 bucket。

\co{ht_bucket} 结构与之前相同，\co{ht_elem} 结构与之前的实现不同之处
仅在于用两元素的链表指针组数组替代了之前的单组链表指针。

在固定大小的 hash table 中，bucket 选择非常直接：
只需将 hash 值转换为对应的 bucket 索引。
相比之下，在调整大小时，更新者还必须确定从旧的和新的 bucket 集中选择哪一个。
如果从旧表中将要选择的 bucket 已经分发到了新表中，
那么应该从新表以及旧表中选择 bucket。
反之，如果从旧表中将要选择的 bucket 尚未分发，
那么应该从旧表中选择 bucket。
\end{fcvref}

\begin{listing}
\input{CodeSamples/datastruct/hash/hash_resize@get_bucket.fcv}
\caption{可调整大小的 Hash-Table Bucket 选择}
\label{lst:datastruct:Resizable Hash-Table Bucket Selection}
\end{listing}

\begin{fcvref}[ln:datastruct:hash_resize:get_bucket]
Bucket 选择如\cref{lst:datastruct:Resizable Hash-Table Bucket Selection}所示，
其中 \co{ht_get_bucket()} 在 \clnrefrange{single:b}{single:e}，
\co{ht_search_bucket()} 在 \clnrefrange{hsb:b}{hsb:e}。
\co{ht_get_bucket()} 函数返回指定 hash table 中与指定 key 对应的 bucket 的引用，
不考虑调整大小的情况。
它还在 \clnref{single:gethash} 上将与 key 对应的 bucket 索引存储到
参数 \co{b} 引用的位置，
并在 \clnref{single:h} 上将与 key 对应的 hash 值存储到
参数 \co{h}（如果非 \co{NULL}）引用的位置。
\Clnref{single:return} 然后返回对应 bucket 的引用。

\co{ht_search_bucket()} 函数在指定的 hash table 版本中搜索指定的 key。
\Clnref{hsb:get_curbkt} 获取与指定 key 对应的 bucket 的引用。
\clnrefrange{hsb:loop:b}{hsb:loop:e} 的循环搜索该 bucket，
因此如果 \clnref{hsb:match} 检测到匹配，
\clnref{hsb:ret_match} 返回指向包含的数据元素的指针。
否则，如果没有匹配，
\clnref{hsb:ret_NULL} 返回 \co{NULL} 表示失败。
\end{fcvref}

\QuickQuiz{
	\cref{lst:datastruct:Resizable Hash-Table Bucket Selection}中的代码
	如何防止调整大小过程推进超过所选 bucket？
}\QuickQuizAnswer{
	它不提供任何这样的保护。
	这是下面描述的更新侧并发控制函数的工作。
}\QuickQuizEnd

\co{ht_get_bucket()} 和 \co{ht_search_bucket()} 的这个实现
允许查找和修改与调整大小操作并发运行。

\begin{listing}
\input{CodeSamples/datastruct/hash/hash_resize@lock_unlock_mod.fcv}
\caption{可调整大小的 Hash-Table 更新侧并发控制}
\label{lst:datastruct:Resizable Hash-Table Update-Side Concurrency Control}
\end{listing}

读侧并发控制由 RCU 提供，如
\cref{lst:datastruct:RCU-Protected Hash-Table Read-Side Concurrency Control}所示，
但更新侧并发控制函数
\co{hashtab_lock_mod()} 和 \co{hashtab_unlock_mod()}
现在必须处理并发调整大小操作的可能性，如
\cref{lst:datastruct:Resizable Hash-Table Update-Side Concurrency Control}所示。

\begin{fcvref}[ln:datastruct:hash_resize:lock_unlock_mod:l]
\co{hashtab_lock_mod()} 在列表的 \clnrefrange{b}{e} 上。
\Clnref{rcu_lock} 进入 RCU 读侧临界区以防止
数据结构在遍历过程中被释放，
\clnref{refhashtbl} 获取当前 hash table 的引用，然后
\clnref{refbucket} 获取该 hash table 中与 key 对应的 bucket 的引用。
\Clnref{acq_bucket} 获取该 bucket 的 lock，这将阻止任何并发调整大小操作
分发该 bucket，尽管如果该 bucket 已经被分发则当然没有效果。
\Clnrefrange{lsp0b}{lsp0e} 将 bucket 指针和指针集索引存储到
\co{ht_lock_state} 结构中各自的字段中，
该结构将信息传递给 \co{hashtab_add()}、\co{hashtab_del()} 和
\co{hashtab_unlock_mod()}。
\Clnref{ifresized} 然后检查并发调整大小操作是否已经将此 bucket 分发到了
新 hash table 中，如果没有，\clnref{lsp1_1} 表示没有已调整大小的 hash bucket，
\clnref{fastret1} 在持有所选 hash bucket 的 lock
（从而阻止并发调整大小操作分发此 bucket）
并且在 RCU 读侧临界区内返回。
通过始终在新表的 lock 之前获取旧表的 lock，
以及利用 RCU 防止在给定时间存在超过两个版本来避免 \IX{Deadlock}，
从而防止死锁（deadlock）环。

否则，并发调整大小操作已经分发了此 bucket，
因此 \clnref{new_hashtbl} 前进到新 hash table，
\clnref{get_newbkt} 选择与 key 对应的 bucket，
\clnref{acq_newbkt} 获取该 bucket 的 lock。
\Clnrefrange{lsp1b}{lsp1e} 将新表的 bucket 指针和
指针集索引存储到
\co{ht_lock_state} 结构的相应字段中，这同样将信息传递给
\co{hashtab_add()}、\co{hashtab_del()} 和 \co{hashtab_unlock_mod()}。
因为此 bucket 已经被调整过大小，而且
\co{hashtab_add()} 和 \co{hashtab_del()} 同时影响新旧两个
\co{ht_bucket} 结构，所以持有两个 lock，分别对应两个 bucket。
此外，\co{ht_lock_state} 结构中每个数组的两个元素都被使用，
其中 \co{[0]} 元素对应旧的 \co{ht_bucket} 结构，
\co{[1]} 元素对应新的结构。
\co{hashtab_lock_mod()} 再次在 RCU 读侧临界区内退出。
\end{fcvref}

\begin{fcvref}[ln:datastruct:hash_resize:lock_unlock_mod:ul]
\co{hashtab_unlock_mod()} 函数释放 \co{hashtab_lock_mod()} 获取的 lock。
\Clnref{relbkt0} 释放旧 \co{ht_bucket} 结构上的 lock。
如果 \clnref{ifbkt1} 确定正在进行调整大小操作（这不太可能），
\clnref{relbkt1} 释放新 \co{ht_bucket} 结构上的 lock。
无论哪种情况，\clnref{rcu_unlock} 退出 RCU 读侧临界区。
\end{fcvref}

\QuickQuiz{
	假设一个线程在调整大小操作期间正在向 hash table 插入元素。
	是什么防止这个插入因后续的调整大小操作在插入完成之前完成而丢失？
}\QuickQuizAnswer{
	第二个调整大小操作将无法越过插入正在进行的 bucket，
	因为插入操作持有 hash table 中一个或两个 hash bucket 上的 lock。
	此外，插入操作在 RCU 读侧临界区内进行。
	正如我们将在检查 \co{hashtab_resize()} 函数时看到的，
	这允许每个调整大小操作使用 \co{synchronize_rcu()} 调用
	来等待插入操作的读侧临界区完成。
	这些 \co{synchronize_rcu()} 调用防止已存在的插入操作
	在调整大小操作之后仍然存在。
}\QuickQuizEnd

\begin{listing}
\input{CodeSamples/datastruct/hash/hash_resize@access.fcv}
\caption{可调整大小的 Hash-Table 访问函数}
\label{lst:datastruct:Resizable Hash-Table Access Functions}
\end{listing}

现在我们已经有了 bucket 选择和并发控制，
可以搜索和更新我们的可调整大小 hash table 了。
\co{hashtab_lookup()}、\co{hashtab_add()} 和 \co{hashtab_del()}
函数如\cref{lst:datastruct:Resizable Hash-Table Access Functions}所示。

\begin{fcvref}[ln:datastruct:hash_resize:access:lkp]
列表中 \clnrefrange{b}{e} 上的 \co{hashtab_lookup()} 函数执行 hash 查找。
\Clnref{get_curtbl} 获取当前 hash table，
\clnref{get_curbkt} 搜索与指定 key 对应的 bucket。
\Clnref{ret} 返回指向搜索元素的指针，
搜索失败时返回 \co{NULL}。
调用者必须在 RCU 读侧临界区内。
\end{fcvref}

\QuickQuiz{
	\cref{lst:datastruct:Resizable Hash-Table Access Functions}中的
	\co{hashtab_lookup()} 函数忽略了并发调整大小操作。
	这不意味着读取者可能会遗漏在调整大小操作期间之前添加的元素吗？
}\QuickQuizAnswer{
	不会。
	我们很快会看到，
	\co{hashtab_add()} 和 \co{hashtab_del()} 函数
	在调整大小操作进行期间保持旧 hash table 是最新的。
}\QuickQuizEnd

\begin{fcvref}[ln:datastruct:hash_resize:access:add]
列表中 \clnrefrange{b}{e} 上的 \co{hashtab_add()} 函数
向 hash table 添加新数据元素。
\Clnref{htbp} 获取要添加新元素的当前 \co{ht_bucket} 结构，
\clnref{i} 获取指针对的索引。
\Clnref{add} 将新元素添加到当前 hash bucket。
如果 \clnref{ifnew} 确定此 bucket 已被分发到新版本的 hash table，
则 \clnref{addnew} 也将新元素添加到对应的新 bucket。
调用者需要处理并发，例如在调用 \co{hashtab_add()} 之前调用
\co{hashtab_lock_mod()}，之后调用 \co{hashtab_unlock_mod()}。
\end{fcvref}

\begin{fcvref}[ln:datastruct:hash_resize:access:del]
列表中 \clnrefrange{b}{e} 上的 \co{hashtab_del()} 函数
从 hash table 中移除已有元素。
\Clnref{i} 获取指针对的索引，
\clnref{del} 从当前表中移除指定元素。
如果 \clnref{ifnew} 确定此 bucket 已被分发到新版本的 hash table，
则 \clnref{delnew} 也从对应的新 bucket 中移除指定元素。
与 \co{hashtab_add()} 一样，调用者负责并发控制，
而这种并发控制足以与并发调整大小操作同步。
\end{fcvref}

\QuickQuiz{
	\cref{lst:datastruct:Resizable Hash-Table Access Functions}中的
	\co{hashtab_add()} 和 \co{hashtab_del()} 函数
	在调整大小操作进行时可能更新两个 hash bucket。
	如果调整大小操作的频率不可忽略，这可能导致性能不佳。
	在这种情况下有没有可能降低更新的成本？
}\QuickQuizAnswer{
	有，至少假设 \co{hashtab_lookup()} 成本的略微增加是可以接受的。
	一种方法如
	\cref{lst:datastruct:Resizable Hash-Table Access Functions (Fewer Updates),%
	lst:datastruct:Resizable Hash-Table Update-Side Locking Function (Fewer Updates)}
	（\path{hash_resize_s.c}）所示。

\begin{listing}
\input{CodeSamples/datastruct/hash/hash_resize_s@access.fcv}
\caption{可调整大小的 Hash-Table 访问函数（更少更新）}
\label{lst:datastruct:Resizable Hash-Table Access Functions (Fewer Updates)}
\end{listing}

\begin{listing}
\input{CodeSamples/datastruct/hash/hash_resize_s@lock_mod.fcv}
\caption{可调整大小的 Hash-Table 更新侧加锁函数（更少更新）}
\label{lst:datastruct:Resizable Hash-Table Update-Side Locking Function (Fewer Updates)}
\end{listing}

	这个版本的 \co{hashtab_add()} 将元素添加到
	旧 bucket（如果尚未调整大小）或新 bucket（如果已调整大小），
	\co{hashtab_del()} 从元素已被插入的任何 bucket 中移除指定元素。
	\co{hashtab_lookup()} 函数在旧 bucket 搜索失败时搜索新 bucket，
	这有一个缺点，即在查找快速路径上增加了开销。
	替代的 \co{hashtab_lock_mod()} 在调整大小操作进行时
	在 \co{->hbp[0]} 和 \co{->hls_idx[0]} 中返回新 bucket 的加锁状态，
	而不是可能更自然的选择 \co{->hbp[1]} 和 \co{->hls_idx[1]}。
	然而，这个不太自然的选择具有简化 \co{hashtab_add()} 的优点。

	对代码的进一步分析留作读者的练习。
}\QuickQuizEnd

\begin{listing*}
\input{CodeSamples/datastruct/hash/hash_resize@resize.fcv}
\caption{可调整大小的 Hash-Table 调整大小}
\label{lst:datastruct:Resizable Hash-Table Resizing}
\end{listing*}

\begin{fcvref}[ln:datastruct:hash_resize:resize]
实际的调整大小由 \co{hashtab_resize()} 执行，如
\cpageref{lst:datastruct:Resizable Hash-Table Resizing}上的
\cref{lst:datastruct:Resizable Hash-Table Resizing}所示。
\Clnref{trylock} 条件性地获取顶层 \co{->ht_lock}，如果获取失败，
\clnref{ret_busy} 返回 \co{-EBUSY} 表示调整大小已在进行中。
否则，\clnref{get_curtbl} 获取当前 hash table 的引用，
\clnrefrange{alloc:b}{alloc:e} 分配所需大小的新 hash table。
如果指定了新的 hash/key 函数集，这些将用于新表，
否则保留旧表的函数。
如果 \clnref{chk_nomem} 检测到内存分配失败，
\clnref{rel_nomem} 释放 \co{->ht_lock}，
\clnref{ret_nomem} 返回失败指示。

\Clnref{get_curidx} 获取当前表的索引，
\clnref{put_curidx} 将其反转值存储到新 hash table，
从而确保两个 hash table 避免覆盖彼此的 linked list。
\Clnref{set_newtbl} 然后通过将新表的引用安装到旧表的 \co{->ht_new} 字段
来开始 bucket 分发过程。
\Clnref{sync_rcu} 确保所有不知道新表的读取者在调整大小操作继续之前完成。

\clnrefrange{loop:b}{loop:e} 循环的每次迭代
将旧 hash table 的一个 bucket 的内容分发到新 hash table。
\Clnref{get_oldcur} 获取旧表当前 bucket 的引用，
\clnref{acq_oldcur} 获取该 bucket 的 spinlock。
\end{fcvref}

\QuickQuiz{
	\begin{fcvref}[ln:datastruct:hash_resize:resize]
	在\cref{lst:datastruct:Resizable Hash-Table Resizing}中的
	\co{hashtab_resize()} 函数中，
	从 \co{hashtab_add()} 和 \co{hashtab_del()} 的角度来看，
	什么保证了 \clnref{set_newtbl} 上对 \co{->ht_new} 的更新
	会被视为发生在 \clnref{update_resize} 上
	对 \co{->ht_resize_cur} 的更新之前？
	换句话说，什么阻止了 \co{hashtab_add()}
	和 \co{hashtab_del()} 解引用从 \co{->ht_new} 加载的 \co{NULL} 指针？
	\end{fcvref}
}\QuickQuizAnswer{
	\begin{fcvref}[ln:datastruct:hash_resize:resize]
	\cref{lst:datastruct:Resizable Hash-Table Resizing}中
	\clnref{sync_rcu} 上的 \co{synchronize_rcu()}
	确保在我们在 \clnref{set_newtbl} 上安装新 hash table 引用
	和在 \clnref{update_resize} 上更新 \co{->ht_resize_cur} 之间，
	所有预先存在的 RCU 读取者已经完成。
	这意味着任何看到 \co{->ht_resize_cur} 非负值的读取者
	不可能在对 \co{->ht_new} 的赋值之前就已经开始，
	因此必须能看到对新 hash table 的引用。

	这就是为什么更新侧的 \co{hashtab_add()} 和 \co{hashtab_del()} 函数
	必须包含在 RCU 读侧临界区中，这由
	\cref{lst:datastruct:Resizable Hash-Table Update-Side Concurrency Control}中的
	\co{hashtab_lock_mod()} 和 \co{hashtab_unlock_mod()} 提供。
	\end{fcvref}
}\QuickQuizEnd

\begin{fcvref}[ln:datastruct:hash_resize:resize]
\clnrefrange{loop_list:b}{loop_list:e} 循环的每次迭代
将一个数据元素从当前旧表 bucket 添加到对应的新表 bucket，
在添加操作期间持有新表 bucket 的 lock。
\Clnref{update_resize} 更新 \co{->ht_resize_cur} 以表示此 bucket 已被分发。
最后，\clnref{rel_oldcur} 释放旧表的 bucket lock。

一旦所有旧表 bucket 都已分发到新表，执行到达 \clnref{rcu_assign}。
\Clnref{rcu_assign} 将新创建的表安装为当前表，
\clnref{sync_rcu_2} 等待所有旧读取者（可能仍在引用旧表）完成。
然后 \clnref{rel_master} 释放调整大小序列化 lock，
\clnref{free} 释放旧 hash table，
最后 \clnref{ret_success} 返回成功。
\end{fcvref}

\QuickQuiz{
	\begin{fcvref}[ln:datastruct:hash_resize:resize]
	为什么在\cref{lst:datastruct:Resizable Hash-Table Resizing}的
	\clnref{update_resize} 上有一个 \co{WRITE_ONCE()}？
	\end{fcvref}
}\QuickQuizAnswer{
	\begin{fcvref}[ln:datastruct:hash_resize:lock_unlock_mod]
	与\cref{lst:datastruct:Resizable Hash-Table Update-Side Concurrency Control}中
	\co{hashtab_lock_mod()} 的 \clnref{l:ifresized} 上的 \co{READ_ONCE()} 一起，
	它告诉编译器对 \co{->ht_resize_cur} 的非初始化访问必须保留，
	因为对 \co{->ht_resize_cur} 的读取确实可能与写入竞争，
	只是不会以改变 ``if'' 条件的方式竞争。
	\end{fcvref}
}\QuickQuizEnd

\subsection{可调整大小的哈希表讨论}
\label{sec:datastruct:Resizable Hash Table Discussion}

\begin{figure}
\centering
\resizebox{2.7in}{!}{\includegraphics{CodeSamples/datastruct/hash/data/hps.resize.2020.09.05a/perftestresize}}
\caption{在 262,144 和 524,288 Bucket 之间调整 Hash Table 大小的开销与总元素数量的关系}
\label{fig:datastruct:Overhead of Resizing Hash Tables Between 262;144 and 524;288 Buckets vs. Total Number of Elements}
\end{figure}

\Cref{fig:datastruct:Overhead of Resizing Hash Tables Between 262;144 and 524;288 Buckets vs. Total Number of Elements}
将可调整大小的 hash table 与其固定大小的对应物进行了比较，
分别在 hash table 中使用 262,144 和 2,097,152 个元素。
该图为每个元素数量显示了三条曲线，一条用于固定大小的 262,144 bucket hash table，
另一条用于固定大小的 524,288 bucket hash table，
第三条用于在 262,144 和 524,288 bucket 之间来回调整的可调整大小 hash table，
每次调整大小操作之间有一毫秒的暂停。

最上面的三条曲线是 262,144 元素 hash table 的。\footnote{
	你只看到两条曲线？
	虚线那条实际上由两条差异很小的曲线组成，
	因此虚线图案看起来不规则。}
虚线曲线对应两个固定大小 hash table，
实线曲线对应可调整大小 hash table。
在这种情况下，短的 hash chain 使得正常查找开销非常低，
以至于调整大小的开销在大部分范围内占主导地位。
特别是，整个 hash table 可以放入 L3 cache。
% Nevertheless, the larger fixed-size hash table has a significant
% performance advantage, so that resizing can be quite beneficial,
% at least given sufficient time between resizing operations: One
% millisecond is clearly too short a time.

下面三条曲线是 2,097,152 元素 hash table 的。
上面的虚线曲线对应 262,144 bucket 的固定大小 hash table，
中间的实线曲线（低 CPU 数量时在中间，高 CPU 数量时在底部）
对应可调整大小 hash table，
另一条曲线对应 524,288 bucket 的固定大小 hash table。
现在每个 bucket 平均有八个元素的事实只能预期会导致性能急剧下降，
实际上图中也显示了这一点。
更糟糕的是，hash table 元素占用 128\,MB，
超出了每个 socket 的 39\,MB L3 cache，
其性能后果类似于\cref{sec:cpu:Costs of Operations}中描述的那些。
由此产生的 cache 溢出意味着即使对于只读基准测试也涉及到内存系统，
从下面三条曲线的亚线性部分可以看出，内存系统可能是一个严重的瓶颈。

\QuickQuiz{
	\cref{fig:datastruct:Overhead of Resizing Hash Tables Between 262;144 and 524;288 Buckets vs. Total Number of Elements}
	中大小 hash table 之间的性能差异，有多少是由于长 hash chain，
	有多少是由于内存系统瓶颈？
}\QuickQuizAnswer{
	回答这个问题的简单方法是用 2,097,152 个元素做另一次运行，
	但这次也使用 2,097,152 个 bucket，
	从而将每个 bucket 的平均元素数降回到一。

\begin{figure}
\centering
\resizebox{2.7in}{!}{\includegraphics{CodeSamples/datastruct/hash/data/hps.resize.2020.09.27a/perftestresizebig}}
\caption{内存系统瓶颈对 Hash Table 的影响}
\label{fig:datastruct:Effect of Memory-System Bottlenecks on Hash Tables}
\end{figure}

	结果由\cref{fig:datastruct:Effect of Memory-System Bottlenecks on Hash Tables}
	中间的三虚线新曲线所示。
	其他六条曲线与
	\cpageref{fig:datastruct:Overhead of Resizing Hash Tables Between 262;144 and 524;288 Buckets vs. Total Number of Elements}上的
	\cref{fig:datastruct:Overhead of Resizing Hash Tables Between 262;144 and 524;288 Buckets vs. Total Number of Elements}
	中的对应曲线相同。
	这条新曲线与下面三条曲线之间的差距粗略衡量了性能差异中
	有多少归因于 hash chain 长度，
	新曲线与上面三条曲线之间的差距粗略衡量了
	有多少归因于内存系统瓶颈。
	新曲线在单个 CPU 时略低于其 262,144 元素的对应曲线，
	表明即使在单个 CPU 上 cache 容量也在略微降低性能。\footnote{
		是的，在硬件架构师看来，cache 是内存系统的一部分。}
	这是可以预期的，因为与其较小的对应物不同，
	2,097,152 bucket 的 hash table 无法放入 L3 cache。
	这条新曲线刚过 28 个 CPU 时上升，这也是可以预期的。
	这一上升是因为第 29 个 CPU 在另一个 socket 上，
	带来了额外的 39\,MB cache 以及额外的内存带宽。

	但大型 hash table 相对于 524,288 bucket（但仍有 2,097,152 个元素）
	的 hash table 的优势随着 CPU 增加而减少，
	这与瓶颈位于内存系统一致。
	在约 400 个 CPU 以上，2,097,152 bucket 的 hash table 实际上
	略微不如 524,288 bucket 的 hash table。
	这不应令人惊讶，因为内存系统是瓶颈，
	更多的 bucket 增加了此工作负载的内存占用。

	细心的读者会注意到上面的"粗略"一词，
	可能对更详细的分析感兴趣。
	欢迎这些读者运行类似的基准测试，
	使用他们可能拥有的任何性能计数器或硬件分析工具。
	这可能是一段漫长而复杂的旅程，但那些勇于踏上它的人
	将获得关于硬件性能及其对软件影响的详细知识。
}\QuickQuizEnd

参考
\cref{tab:cpu:CPU 0 View of Synchronization Mechanisms on 8-Socket System With Intel Xeon Platinum 8176 CPUs at 2.10GHz}
的最后一列，我们回忆前 28 个 CPU 在第一个 socket 中，
每核一个 CPU，这解释了可调整大小 hash table 在 28 个 CPU 之后
性能急剧下降的原因。
尽管这种下降很急剧，请记住它是由于不断来回调整大小造成的。
显然一次性调整到 524,288 个 bucket 会更好，
或者更好的是，做一次八倍的调整到 2,097,152 个元素，
从而将每个 bucket 的平均元素数降低到产生上面三条曲线的运行所享有的水平。

从这些数据得出的关键点是，RCU 保护的可调整大小 hash table
的性能和可扩展性几乎与其固定大小的对应物一样好。
实际调整大小操作期间的性能当然会因更新每个元素指针导致的 cache miss
而有所下降，当内存系统成为瓶颈时这种效应最为显著。
这表明 hash table 应该大幅度调整大小，
并且应该应用滞后机制以防止过于频繁的调整大小操作导致性能下降。
在内存充裕的环境中，hash table 大小的增加
应该比减少激进得多。

另一个关键点是，虽然 \co{hashtab} 结构是不可分区的，
但它也是以读为主的，这暗示应使用 RCU\@。
鉴于这种可调整大小 hash table 的性能和可扩展性
非常接近 RCU 保护的固定大小 hash table，
我们必须得出结论，这种方法相当成功。

最后，需要注意的是插入、删除和查找可以与调整大小操作并发进行。
这种并发性在调整大型 hash table 大小时至关重要，
特别是对于必须满足严格响应时间约束的应用。

当然，\co{ht_elem} 结构的一对指针集确实带来了一些内存开销，
这将在下一节讨论。

\subsection{其他可调整大小的哈希表}
\label{sec:datastruct:Other Resizable Hash Tables}

本节前面描述的可调整大小 hash table 的一个缺点是内存消耗。
每个数据元素有两对 linked list 指针而不是一对。
是否有可能创建一个只需一对指针的 RCU 保护可调整大小 hash table？

事实证明答案是肯定的。
Josh Triplett 等人~\cite{Triplett:2011:RPHash}
开发了一种\emph{相对论 hash table}，它增量地分割和合并对应的 hash chain，
使得读取者在调整大小操作的整个过程中始终看到有效的 hash chain。
这种增量分割和合并依赖于这样一个事实：
读取者看到应该在其他 hash chain 中的数据元素是无害的——
当这种情况发生时，读取者会因为 key 不匹配而简单地忽略多余的数据元素。

\begin{figure}
\centering
\resizebox{3in}{!}{\includegraphics{datastruct/zipperhashshrink}}
\caption{缩小相对论 Hash Table}
\label{fig:datastruct:Shrinking a Relativistic Hash Table}
\end{figure}

将相对论 hash table 缩小两倍的过程如
\cref{fig:datastruct:Shrinking a Relativistic Hash Table}所示，
在这种情况下是将一个双 bucket hash table 缩小为单 bucket hash table，
也称为线性链表。
这个过程通过将旧的较大 hash table 中的成对 bucket 合并为
新的较小 hash table 中的单个 bucket 来工作。
为使这个过程正确工作，我们显然需要约束两个表的 hash 函数。
一种约束是对两个表使用相同的底层 hash 函数，
但在从大缩小到小时丢弃低位 bit。
例如，旧的双 bucket hash table 将使用值的前两个 bit，
而新的单 bucket hash table 可以使用值的最高位 bit。
这样，旧的大 hash table 中相邻的偶数和奇数 bucket 对
就可以合并为新的小 hash table 中的单个 bucket，
同时仍然有一个 hash 值覆盖该单个 bucket 中的所有元素。

初始状态显示在图的顶部，时间从上到下前进，从初始状态~(a) 开始。
缩小过程从分配新的较小 bucket 数组开始，
让这个新的较小数组的每个 bucket 引用旧大 hash table 中
对应对中一个 bucket 的第一个元素，产生状态~(b)。

然后两个 hash chain 被链接在一起，产生状态~(c)。
在这个状态下，查找偶数编号元素的读取者看不到变化，
查找元素~1 和~3 的读取者同样看不到变化。
然而，查找其他奇数的读取者也会遍历元素~0 和~2。
这是无害的，因为任何奇数都会与这两个元素比较不等。
虽然有一些性能损失，但另一方面，这正是新的小 hash table
完全就位后将会经历的同样性能损失。

接下来，新的小 hash table 被开放给读取者访问，产生状态~(d)。
注意，较早的读取者可能仍在遍历旧的大 hash table，
因此在这个状态下两个 hash table 都在使用中。

下一步是等待所有预先存在的读取者完成，产生状态~(e)。
在这个状态下，所有读取者都在使用新的小 hash table，
因此旧的大 hash table 的 bucket 可以被释放，产生最终状态~(f)。

\begin{figure}
\centering
\resizebox{3in}{!}{\includegraphics{datastruct/zipperhashgrow}}
\caption{增长相对论 Hash Table}
\label{fig:datastruct:Growing a Relativistic Hash Table}
\end{figure}

增长相对论 hash table 是缩小过程的逆过程，
但需要更多的 grace period 步骤，如
\cref{fig:datastruct:Growing a Relativistic Hash Table}所示。
初始状态~(a) 在该图的顶部，时间从上到下前进。

我们从分配新的大型双 bucket hash table 开始，产生状态~(b)。
注意每个新 bucket 引用该 bucket 的第一个目标元素。
这些新 bucket 发布给读取者，产生状态~(c)。
经过 grace period 操作，所有读取者都在使用新的大型 hash table，
产生状态~(d)。
在这个状态下，只有遍历偶数值 hash bucket 的读取者才会遍历元素~0，
因此它现在被涂成白色。

此时，旧的小 hash bucket 可以被释放，尽管许多实现使用这些旧 bucket
来跟踪将项目列表"解压"到各自新 bucket 中的进度。
第一组连续偶数编号元素中的最后一个偶数编号元素的
指向下一个的指针现在更新为引用随后的偶数编号元素。
经过随后的 grace period 操作，结果是状态~(e)。
竖向箭头指示下一个要解压的元素，
元素~1 现在被涂成黑色，表示只有遍历奇数值 hash bucket 的读取者才能到达它。

接下来，第一组连续奇数编号元素中的最后一个奇数编号元素的
指向下一个的指针现在更新为引用随后的奇数编号元素。
经过随后的 grace period 操作，结果是状态~(f)。
最后一次解拉链操作（包括一次 grace period 操作）产生最终状态~(g)。

简而言之，相对论 hash table 以调整大小期间产生的额外 grace period 为代价，
减少了每个元素的链表指针数量。
这些额外的 grace period 通常不是问题，
因为插入、删除和查找可以与调整大小操作并发进行。

然而，当 hash table 快速增长时，大量的 grace period 可能会
延迟下一次调整大小操作的开始。
减少 grace period 数量的一种方法是同时解压所有 bucket，
这样 grace period 的数量将受限于单个 bucket 中最大元素数量，
而不是 hash table 中的总元素数量。
另一种方法，同样由 Josh Triplett 提出，是维护一个单独的指针
指向当前正在从旧 bucket 移动到新 bucket 的元素。
读取者然后按顺序检查旧 bucket、单独的指针和新 bucket。
这需要读侧内存排序，会减慢读取者，但它大大减少了
调整大小时所需的 grace period 数量。
Linux 内核中使用了这种方法。

事实证明，可以将每个元素的内存开销从一对指针减少到单个指针，
同时仍然保留 $\O{1}$ 的删除操作。
这是通过用 RCU 保护增强 split-order
list~\cite{OriShalev2006SplitOrderListHash}
来实现的~\cite{MathieuDesnoyers2009URCU,PaulMcKenney2013LWNURCUhash}。
Hash table 中的数据元素被排列成一个单一的有序 linked list，
每个 hash bucket 引用该 bucket 中的第一个元素。
元素通过在其指向下一个的指针字段中设置低位 bit 来删除，
这些元素由后续遇到它们的遍历从链表中移除。

这种 RCU 保护的 split-order list 很复杂，
但为所有插入、删除和查找操作提供了 lock-free 进展保证。
这种保证对于实时应用可能很重要。
然而，一个缺点是 key 在查找时必须进行位反转，这会减慢读取者。
但对于不在意这种减速的人，
可以从用户空间 RCU 库的最新版本获得一个实现~\cite{MathieuDesnoyers2009URCU}。

\section{其他数据结构}
\label{sec:datastruct:Other Data Structures}
%
\epigraph{所有的生活都是实验。
	  你做的实验越多越好。}
	 {Ralph Waldo Emerson}

前面几节聚焦于由于可分区性
（\cref{sec:datastruct:Partitionable Data Structures}）、
高效处理以读为主的访问模式
（\cref{sec:datastruct:Read-Mostly Data Structures}）、
或应用以读为主技术来避免不可分区性
（\cref{sec:datastruct:Non-Partitionable Data Structures}）
而增强并发性的数据结构。
本节简要回顾其他数据结构。

Hash table 用于并行的最大优势之一是它完全可分区，至少在不调整大小时是这样。
保持可分区性和大小独立性的一种方法是使用 radix tree，也称为 trie。
Trie 对搜索 key 进行分区，使用每个连续的 key 分区来遍历 trie 的下一层。
因此，trie 可以被看作是一组嵌套的 hash table，从而提供所需的可分区性。
Trie 的一个缺点是稀疏的 key 空间可能导致内存使用效率低下。
有多种压缩技术可以用来解决这个缺点，
包括在遍历之前将 key 值 hash 到更小的 key 空间~\cite{RobertOlsson2007Trash}。
Radix tree 在实践中被大量使用，包括在 Linux
内核中~\cite{NickPiggin2006radixtree}。

Hash table 和 trie 的一个重要特例是也许最古老的数据结构——
数组及其多维对应物矩阵。
矩阵的完全可分区特性在并发数值算法中被大量利用。

自平衡树在顺序代码中被大量使用，
AVL 树和红黑树可能是最知名的例子~\cite{ThomasHCorman2001Algorithms}。
早期并行化 AVL 树的尝试复杂且不一定高效~\cite{Ellis80}，
然而，红黑树的更近期工作通过对读取者使用 RCU 和
对更新使用 hash lock 数组\footnote{
	以 swissTM~\cite{AleksandarDragovejic2011STMnotToy} 的形式，
	这是软件事务内存（transactional memory）的一种变体，
	其中开发者标记非共享访问。}
来保护读取和更新，提供了更好的性能和
可扩展性~\cite{PhilHoward2011RCUTMRBTree,PhilipWHoward2013RCUrbtree}。
事实证明红黑树激进地重新平衡，这对顺序程序效果好，
但对并行使用不一定那么好。
因此，最近的工作使用了 RCU 保护的"盆景树"，
它不那么激进地重新平衡~\cite{AustinClements2012RCULinux:mmapsem}，
以牺牲最优树深度来获得更高效的并发更新。

并发 skip list 非常适合 RCU 读取者，
实际上代表了一种类似 RCU 技术的早期学术使用~\cite{Pugh90}。
% @@@ Reference skiplist chapter once added.

并发双端队列在\cref{sec:SMPdesign:Double-Ended Queue}中讨论过，
并发栈和队列有很长的历史~\cite{Treiber86}，
尽管通常不会有最令人印象深刻的性能或可扩展性。
尽管如此，它们仍然是并发库的常见
特性~\cite{PaulMcKenney2013LWNURCUqueuestack}。
研究人员最近提出放松栈和队列的排序
约束~\cite{Shavit:2011:DSM:1897852.1897873}，
一些工作表明放松排序的队列实际上比严格的 FIFO
队列具有更好的排序
属性~\cite{AndreasHaas2012FIFOisnt,ChristophMKirsch2012FIFOisntTR,AndreasHaas2013CFRelaxedQueues}。

继续研究并发数据结构很可能会产生具有令人惊讶属性的新算法。

\section{总结}
\label{sec:datastruct:Summary}
%
\epigraph{只有一件事比从经验中学习更痛苦，那就是不从经验中学习。}
	 {Archibald MacLeish}

本章主要聚焦于 hash table，包括不完全可分区的可调整大小 hash table。
\Cref{sec:datastruct:Other Data Structures}简要概述了几种非 hash table 数据结构。
尽管如此，这些对 hash table 的阐述是了解围绕高性能可扩展数据访问的
许多问题的优秀入门，包括：

\begin{enumerate}
\item	完全分区的数据结构在小型系统（例如单 socket 系统）上工作良好。
\item	更大的系统除了完全分区外还需要访问局部性。
\item	以读为主的技术（如 hazard pointer 和 RCU）为以读为主的工作负载
	提供良好的访问局部性，因此即使在更大的系统上也能提供优秀的性能和可扩展性。
\item	以读为主的技术也适用于某些类型的不可分区数据结构，
	如可调整大小的 hash table。
\item	大型数据结构可能溢出 CPU cache，降低性能和可扩展性。
\item	通过将数据结构专门化为特定工作负载，可以获得额外的性能和可扩展性，
	例如用 32 位整数替代通用 key。
\item	虽然可移植性和极致性能的需求经常冲突，
	但有一些数据结构布局技术可以在这两组需求之间取得良好的平衡。
\end{enumerate}

话虽如此，没有可靠性的性能和可扩展性用处不大，
因此下一章将涵盖验证。

\QuickQuizAnswersChp{qqzdatastruct}
